\documentclass[11pt]{article}

\usepackage{amsmath,amssymb,amsthm}
\usepackage[margin=1in]{geometry}

\newtheorem{theorem}{Theorem}
\newtheorem{lemma}{Lemma}

\DeclareMathOperator{\Bin}{Bin}

\begin{document}
\title{Submission C solution to Problem 3}
\date{}
\maketitle


\begin{abstract}
We determine all \(p\in[0,1]\) such that, for every finite probability vector
\((w_1,\ldots,w_m)\) and i.i.d. Bernoulli\((p)\) random variables
\(v_1,\ldots,v_m\), one has
\[
\mathbb P\!\left(\sum_{i=1}^m w_i v_i\ge p\right)\ge p.
\]
The answer is
\[
\boxed{\,p\in[0,1/3]\cup\{1/2,1\}\, }.
\]
\end{abstract}

\section*{External inequalities used}

We use two standard sharp inequalities.

\begin{lemma}[Hoeffding--Pinelis comparison for weighted Bernoulli sums]
\label{lem:hoeffding}
Let \(0<\theta<1\), let \(\xi_1,\ldots,\xi_m\) be i.i.d. Bernoulli\((\theta)\), and let
\(a_i\ge0\), \(\sum_i a_i=1\). Then for every real \(x\),
\[
\mathbb P\!\left(\sum_{i=1}^m a_i\xi_i<x\right)
\le
\sup_{n\ge1}
\mathbb P\!\left(\frac1n B_{n,\theta}<x\right),
\]
where \(B_{n,\theta}\sim \Bin(n,\theta)\).
\end{lemma}

This is the Bernoulli specialization of Hoeffding's extremal theorem for sums of
bounded independent random variables; see Hoeffding \cite[Theorem 4]{Hoeffding1956}.
For the left-tail formulation in essentially the above form, see also Pinelis
\cite[Corollary 3.7]{Pinelis2016}.

\begin{lemma}[Ramanujan--Jogdeo--Samuels binomial inequality]
\label{lem:binomial}
If \(0\le \theta\le 1/3\) and \(B_{n,\theta}\sim \Bin(n,\theta)\), then
\[
\mathbb P(B_{n,\theta}<n\theta)\le 1-\theta
\]
for every integer \(n\ge1\).
\end{lemma}

Equivalently,
\[
\mathbb P(B_{n,\theta}\ge n\theta)\ge \theta .
\]
This is a classical binomial-tail inequality proved by Jogdeo and Samuels
\cite{JogdeoSamuels1968}; it is often discussed as a Ramanujan-type binomial
inequality.

\section*{Main theorem}

\begin{theorem}
\label{thm:main}
Let \(p\in[0,1]\). The following assertion holds for every \(m\ge1\), every
choice of nonnegative weights \(w_1,\ldots,w_m\) with \(\sum_i w_i=1\), and
i.i.d. Bernoulli\((p)\) random variables \(v_1,\ldots,v_m\):
\[
\mathbb P\!\left(\sum_{i=1}^m w_i v_i\ge p\right)\ge p.
\]
This is true if and only if
\[
p\in[0,1/3]\cup\{1/2,1\}.
\]
\end{theorem}

\begin{proof}
We prove sufficiency and necessity separately.

\medskip

\noindent\textbf{Sufficiency for \(0\le p\le 1/3\).}
The case \(p=0\) is trivial, since the event
\(\sum_i w_i v_i\ge0\) has probability \(1\). Assume \(0<p\le1/3\), and set
\[
S=\sum_{i=1}^m w_i v_i .
\]
By Lemma~\ref{lem:hoeffding}, with \(\theta=p\) and \(x=p\),
\[
\mathbb P(S<p)
\le
\sup_{n\ge1}\mathbb P(B_{n,p}<np).
\]
By Lemma~\ref{lem:binomial}, each term in the supremum is at most \(1-p\).
Therefore
\[
\mathbb P(S<p)\le 1-p,
\]
and hence
\[
\mathbb P(S\ge p)=1-\mathbb P(S<p)\ge p.
\]

\medskip

\noindent\textbf{Sufficiency for \(p=1/2\).}
Let
\[
S=\sum_{i=1}^m w_i v_i .
\]
Since \(v_i\sim\operatorname{Bernoulli}(1/2)\), the vector
\((1-v_1,\ldots,1-v_m)\) has the same distribution as
\((v_1,\ldots,v_m)\). Thus \(1-S\) has the same distribution as \(S\). Hence
\[
\mathbb P(S>1/2)=\mathbb P(S<1/2).
\]
Consequently,
\[
\mathbb P(S\ge1/2)
=
\mathbb P(S>1/2)+\mathbb P(S=1/2)
=
\frac{1+\mathbb P(S=1/2)}2
\ge \frac12.
\]

\medskip

\noindent\textbf{Sufficiency for \(p=1\).}
If \(p=1\), then \(v_i=1\) almost surely for every \(i\), so
\[
\sum_i w_i v_i=\sum_i w_i=1
\]
almost surely. Thus
\[
\mathbb P\!\left(\sum_i w_i v_i\ge1\right)=1.
\]

\medskip

\noindent\textbf{Necessity: failure for \(1/3<p<1/2\).}
Take \(m=3\) and \(w_1=w_2=w_3=1/3\). Let
\[
K=v_1+v_2+v_3\sim \Bin(3,p).
\]
For \(1/3<p<1/2\),
\[
\frac K3\ge p
\quad\Longleftrightarrow\quad
K\ge2.
\]
Therefore
\[
\mathbb P\!\left(\frac K3\ge p\right)
=
\mathbb P(K\ge2)
=
3p^2(1-p)+p^3
=
p^2(3-2p).
\]
Now
\[
p^2(3-2p)<p
\]
is equivalent to
\[
p(3-2p)<1,
\]
or
\[
2p^2-3p+1>0,
\]
i.e.
\[
(2p-1)(p-1)>0.
\]
This holds for \(1/3<p<1/2\). Hence the desired inequality fails throughout
\((1/3,1/2)\).

\medskip

\noindent\textbf{Necessity: failure for \(1/2<p<1\).}
Take \(m=2\) and \(w_1=w_2=1/2\). Let
\[
K=v_1+v_2\sim \Bin(2,p).
\]
For \(p>1/2\),
\[
\frac K2\ge p
\quad\Longleftrightarrow\quad
K=2.
\]
Thus
\[
\mathbb P\!\left(\frac K2\ge p\right)=p^2<p
\]
for every \(p\in(1/2,1)\). Hence the desired inequality fails throughout
\((1/2,1)\).

Combining the sufficiency and necessity parts proves the theorem.
\end{proof}

\begin{thebibliography}{9}

\bibitem{Hoeffding1956}
W.~Hoeffding,
\newblock On the distribution of the number of successes in independent trials,
\newblock \emph{Annals of Mathematical Statistics} \textbf{27} (1956), 713--721.

\bibitem{JogdeoSamuels1968}
K.~Jogdeo and S.~M. Samuels,
\newblock Monotone convergence of binomial probabilities and a generalization of Ramanujan's equation,
\newblock \emph{Annals of Mathematical Statistics} \textbf{39} (1968), 1191--1195.

\bibitem{Pinelis2016}
I.~Pinelis,
\newblock Optimal binomial, Poisson, and normal left-tail domination for sums of nonnegative random variables,
\newblock \emph{Electronic Journal of Probability} \textbf{21} (2016), article no.~47.

\end{thebibliography}

\end{document}